% ============================================================
% CHAPTER: PROJECT EXECUTION
% ============================================================

\chapter{Project Execution}

\section{Problem Understanding and Requirement Analysis}

\begin{enumerate}
    \item \textbf{Define the project scope:}
    \begin{itemize}
        \item Understand the goal of detecting emotions and depression from vocal data.
        \item Clearly identify the emotions to be detected (e.g., happy, sad, angry, etc.) and depression levels (binary or continuous).
    \end{itemize}
    \item \textbf{Study the datasets:}
    \begin{itemize}
        \item Analyze the RAVDESS (emotional speech and song) and DAIC (depression-related interviews) datasets.
        \item Understand the features available (e.g., audio clips, labels) and their structure (e.g., file formats, sample rate).
        \item Review the annotations (emotion labels for RAVDESS and depression scores for DAIC).
    \end{itemize}
    \item \textbf{Set project requirements:}
    \begin{itemize}
        \item Identify the computational resources required (e.g., hardware, memory, GPU).
        \item Determine software tools (e.g., Python, TensorFlow, PyTorch).
        \item Set performance metrics for evaluation (e.g., accuracy, F1-score, recall, precision).
    \end{itemize}
\end{enumerate}

\section{Dataset Preparation}

\begin{enumerate}
    \item \textbf{Preprocess the RAVDESS dataset for emotion classification:}
    \begin{itemize}
        \item Extract features using audio feature extraction techniques like Mel-Frequency Cepstral Coefficients (MFCCs).
        \item Optionally extract other features like spectral features, pitch, and energy.
        \item Normalize and label data by rescaling amplitudes and categorizing according to emotional categories (happy, sad, angry, etc.).
    \end{itemize}
    \item \textbf{Preprocess the DAIC dataset for depression detection:}
    \begin{itemize}
        \item Extract features like speech duration, tone, pitch, pause duration, and other relevant vocal characteristics.
        \item Annotate with depression labels based on depression severity (e.g., binary: depressed/non-depressed, or continuous scores).
        \item Align the annotations with the features extracted from the audio data.
    \end{itemize}
    \item \textbf{Split datasets:}
    \begin{itemize}
        \item Split both RAVDESS and DAIC datasets into training, validation, and test sets (80\% training, 20\% testing).
        \item Ensure that the data is balanced for both emotion and depression labels.
    \end{itemize}
\end{enumerate}

\section{Model Development}

\begin{enumerate}
    \item \textbf{Develop CNN model for emotion detection:}
    \begin{itemize}
        \item Design architecture with layers like convolutional layers, pooling layers, and fully connected layers to classify emotions.
        \item Consider using 1D convolutions for sequence-based audio features or 2D convolutions if using spectrograms.
    \end{itemize}
    \item \textbf{Hyperparameter tuning:}
    \begin{itemize}
        \item Select appropriate hyperparameters like learning rate, batch size, and number of layers.
        \item Use techniques such as grid search or random search for hyperparameter optimization.
    \end{itemize}
    \item \textbf{Develop CNN model for depression analysis:}
    \begin{itemize}
        \item Design a second CNN model to classify or regress depression levels.
        \item Incorporate vocal features like tone, pitch, and speech duration.
        \item Account for subtle features to ensure the model captures subtle vocal cues indicative of depression, such as pauses and tone changes.
        \item Apply data augmentation techniques (e.g., pitch shifting, adding noise) to improve the robustness of the model.
    \end{itemize}
\end{enumerate}

\section{Training and Optimization}

\begin{enumerate}
    \item \textbf{Train CNN models:}
    \begin{itemize}
        \item Monitor training metrics including training accuracy, loss, precision, recall, and F1-score throughout the training process.
        \item Use validation sets to monitor performance and avoid overfitting.
    \end{itemize}
    \item \textbf{Use callbacks:}
    \begin{itemize}
        \item Implement early stopping to halt training if performance on the validation set stops improving.
        \item Optionally use model checkpoints to save the best performing model during training.
    \end{itemize}
    \item \textbf{Validate models:}
    \begin{itemize}
        \item Fine-tune hyperparameters based on the performance on the validation dataset.
        \item Optionally use cross-validation to ensure that the model generalizes well across different subsets of the data.
    \end{itemize}
    \item \textbf{Test the final models:}
    \begin{itemize}
        \item Evaluate the final models on the held-out test data and report the performance metrics (e.g., accuracy, F1-score).
        \item Compare the performance with baseline models or existing approaches to validate the effectiveness of the CNN models.
    \end{itemize}
\end{enumerate}

\section{Integration with User Interface}

\begin{enumerate}
    \item \textbf{Develop the UI:}
    \begin{itemize}
        \item Create a graphical user interface (UI) using JavaScript, HTML, CSS (for local use) or Flask/Django.
    \end{itemize}
    \item \textbf{Display results:}
    \begin{itemize}
        \item After processing the audio, display results for both emotion classification and depression detection clearly in the UI.
    \end{itemize}
    \item \textbf{Integrate CNN models with the UI:}
    \begin{itemize}
        \item Ensure smooth communication between the backend (CNN models) and the frontend (UI).
        \item Optimize for real-time analysis to ensure that predictions are fast and can be performed in real-time or near real-time to provide an interactive experience.
    \end{itemize}
\end{enumerate}

\section{Evaluation and Validation}

\begin{enumerate}
    \item \textbf{Conduct thorough testing:}
    \begin{itemize}
        \item Test the system with real-world audio samples to evaluate its robustness and reliability.
    \end{itemize}
    \item \textbf{Gather user feedback:}
    \begin{itemize}
        \item Collect feedback from users to assess the system's usability and effectiveness in real-world scenarios.
    \end{itemize}
    \item \textbf{Compare with existing methods:}
    \begin{itemize}
        \item Benchmark the system's performance against existing emotion detection and depression analysis methods to assess improvements.
    \end{itemize}
\end{enumerate}

\section{Future Enhancements}

\begin{enumerate}
    \item \textbf{Expand with additional datasets:}
    \begin{itemize}
        \item Incorporate additional emotion and depression-related datasets to enhance model accuracy and generalizability.
    \end{itemize}
    \item \textbf{Explore alternative models:}
    \begin{itemize}
        \item Investigate the use of other machine learning models like RNNs or transformers for improved performance.
    \end{itemize}
    \item \textbf{Enhance system functionality:}
    \begin{itemize}
        \item Integrate mental health interventions or suggestions based on the detected emotion or depression level (e.g., links to resources or support).
    \end{itemize}
\end{enumerate}
